\documentclass[11pt]{article}
\usepackage[margin=1in]{geometry}
\usepackage{times}
\usepackage{amsmath,amssymb}
\usepackage{graphicx}
\usepackage{booktabs}
\usepackage{hyperref}
\usepackage{authblk}
\usepackage[T1]{fontenc}
\hypersetup{colorlinks=true,linkcolor=blue,citecolor=blue,urlcolor=blue}

\title{A Quiet Failure in Calibrated Virtual Screening: Marginal Conformal Prediction Under-Covers the Minority Class, and a Class-Conditional Fix Recovers It}
\author[1]{Muhammadjon Tursunbadalov}
\author[1]{Mustafojon Tursunbadalov}
\affil[1]{School of Science and Technology, Champions College Prep, United States}
\date{}

\begin{document}
\maketitle
\noindent\texttt{\{yaiuftuto, yunmt123\}@gmail.com} \\
\noindent Corresponding author: Muhammadjon Tursunbadalov (yaiuftuto@gmail.com)

\begin{abstract}
Conformal prediction is being adopted across drug discovery as a way to put an honest number on
model reliability. The appeal is its guarantee: pick an error rate $\alpha$, and the method
returns prediction sets that contain the true label with probability at least $1-\alpha$. We show
that taking this guarantee at face value can be dangerous on the imbalanced datasets that
screening actually deals with. Across five binary classification datasets, including a real
high-throughput bioactivity screen and an external clinical-toxicity dataset, standard (marginal)
conformal prediction hits its global 90\% coverage target while leaving the minority class badly
exposed---realized minority-class coverage falls to 59.2\% on the BBBP blood--brain-barrier
penetration dataset, to 40.5\% on the Tox21 SR-ARE stress-response dataset, to 2.5\% on the
ClinTox clinical-trial-toxicity dataset, and to 2.1\% on a real PubChem bioactivity screen (AID
651631, p53 pathway), where the rare class is all but abandoned. The failure is not tied to one
model. A random forest on circular fingerprints, a graph convolutional network that reads
molecular structure directly, and a frozen pretrained chemical language model all reproduce it
($p<0.001$ over ten random splits in every case), and the size of the gap tracks how
well-calibrated a model already is on rare labels rather than anything about its architecture.
Rebalancing the classifier itself (class-weighted training) does not fix the failure, which
localizes the problem to the calibration step rather than the model. We explain the effect with a
conservation identity: because overall coverage is fixed, the minority's shortfall equals the
majority's surplus amplified by the imbalance ratio, a relation that predicts the measured BBBP gap
to within about a point and orders the severity across datasets. We extend this identity into a
bounded, validated pre-flight diagnostic that predicts minority coverage to within 2.4 percentage
points of its measured value across a thirteen-fold range of imbalance severity, before a full
calibration study is run. The failure survives realistic scaffold splits, and aggregate accuracy
and overall coverage stay reassuringly high throughout, which is exactly why it is easy to miss.
Class-conditional (Mondrian) conformal prediction closes the gap on every dataset, returning
minority coverage to target for a modest increase in prediction-set size; adaptive prediction sets
(APS), by contrast, reach nominal minority coverage only by returning uninformative two-label sets
for nearly every minority compound. We localize the marginal failures to a handful of generic
molecular scaffolds---plain benzene and pyridine cores that occur in both classes---propose a
one-number diagnostic (the minority coverage gap), and show with a direction-aware cost model that
abstaining on the affected compounds flips a screening campaign from net-negative to net-positive
utility, with the size of that swing set by whether the campaign's goal is to find or to avoid the
class in question. The marginal--conditional distinction is well known in the conformal
literature; our contribution is to demonstrate on real chemistry how severe and how invisible it
becomes under imbalance, to explain it quantitatively, to turn that explanation into a validated
pre-flight predictive check, to show ordinary class rebalancing does not substitute for
class-conditional calibration, and to lay out a practical, cost-aware protocol that restores
per-class reliability.
\end{abstract}

\noindent\textbf{Keywords:} conformal prediction, uncertainty quantification, virtual screening,
class imbalance, selective prediction, molecular property prediction

\section{Introduction}

A screening model is only useful if you can trust the cases where it says it is sure. Modern
classifiers for molecular property prediction reach high aggregate accuracy, but a single accuracy
number says nothing about which predictions deserve confidence. That gap has pushed conformal
prediction into cheminformatics. Conformal methods wrap an existing classifier and return, for
each compound, a set of labels that provably contains the truth at a user-chosen rate
\cite{vovk2005,angelopoulos2023}. The guarantee is distribution-free and finite-sample. It needs no
assumption about the model being correct, only that calibration and test compounds are
exchangeable. For a field that routinely commits expensive wet-lab resources on the strength of a
model's ranking, a guarantee like that is worth a great deal.

There is a catch that is easy to read past. The standard guarantee is marginal: coverage holds on
average over the whole test distribution. It does not promise anything about any particular
subgroup. On a balanced problem the distinction rarely bites. Screening problems are almost never
balanced. The interesting label---a toxic compound, a non-penetrant candidate, an active
hit---is usually the rare one, and the rare one is the whole reason for running the campaign. A
method that keeps its overall promise by being nearly perfect on the common class and quietly poor
on the rare class is, for screening, keeping the wrong promise.

This paper measures how badly that plays out, and shows it is not a corner case. On four imbalanced
tasks, including one drawn from a real high-throughput bioactivity screen, marginal conformal
prediction meets its 90\% target overall while delivering minority-class coverage as low as 2.1\%.
We then ask whether the effect is an artifact of one model family. It is not. The same collapse
appears with fingerprint-based random forests, with a graph neural network, and with a pretrained
transformer that treats a molecule as a string of text---three ways of representing a molecule that
share almost nothing mechanically. The severity of the collapse turns out to depend on how
well-calibrated the base model already is on rare labels, not on which representation it uses. That
is a statement about the calibration step, not about deep learning. We further ask whether the
failure is simply a symptom of an unweighted classifier, and whether a more modern conformal score
(adaptive prediction sets) avoids it without resorting to class-conditional calibration. Neither
turns out to be true.

What makes the failure dangerous in practice is that it is invisible from the top. Overall
coverage, overall accuracy, ROC-AUC---all of the dashboard numbers a team would normally
check---stay healthy while minority coverage falls through the floor. A practitioner who trusts
the conformal guarantee, as advertised, would have no reason to look closer.

The fix is not new and we do not claim it as ours. Class-conditional, or Mondrian, conformal
prediction calibrates a separate threshold per class \cite{vovk2013}, which restores coverage
within each group by construction. The marginal-versus-conditional gap is textbook conformal
theory \cite{angelopoulos2023}. What we contribute is the empirical and explanatory case that this
textbook gap is a live hazard in virtual screening: we quantify it on real datasets including a
genuine bioactivity screen, show it is architecture-agnostic, explain its size with a simple
identity, turn that identity into a checkable pre-flight diagnostic, show it is not fixed by
rebalancing the classifier or by switching to a more modern conformal score, trace it to specific
chemical substructures, and tie it to a decision-cost model that makes the practical stakes
concrete. Our aim is less ``here is a method'' and more ``here is a vulnerability, here is why it
happens, here is how to see it coming, and here is how to close it.''

Concretely, the paper makes seven claims, each backed by experiment:

\begin{enumerate}
\item On imbalanced molecular datasets, marginal conformal prediction satisfies its global
coverage target while under-covering the minority class, and the shortfall grows with imbalance,
from a near-zero gap on near-balanced BACE to near-total collapse on heavily imbalanced ClinTox
and on a real PubChem bioactivity screen.
\item The collapse reproduces across three unrelated model architectures at $p<0.001$ each, with
severity governed by baseline minority calibration rather than architecture.
\item A conservation identity explains the effect: the minority shortfall equals the majority
surplus amplified by the imbalance ratio, which predicts the measured gaps and their ordering
across datasets. We extend this identity into a bounded, validated \emph{pre-flight} diagnostic
that predicts minority coverage from prevalence and one inexpensive calibration pass, before a
full study is run, to within 2.4 percentage points across five datasets and a thirteen-fold range
of imbalance.
\item Standard aggregate metrics conceal the failure, which we demonstrate with a
selective-prediction protocol and a direction-aware cost analysis, and the failure persists under
realistic scaffold splits.
\item Rebalancing the base classifier via class-weighted training does not fix the coverage
failure, which is further evidence the failure lives in the calibration step, not the model.
\item Adaptive prediction sets (APS), a more modern conformal score, reach nominal minority
coverage but only by returning uninformative two-label sets for nearly every minority compound, so
neither default marginal method is usable for screening decisions under imbalance.
\item The confidently wrong minority predictions concentrate on a small set of common,
low-information scaffolds, which we present as a chemical case study.
\end{enumerate}

\section{Related Work}

\textbf{Conformal prediction and its conditional variants.} The framework originates with Vovk and
colleagues \cite{vovk2005}, who established the exchangeability-based coverage guarantee. The
accessible modern treatment by Angelopoulos and Bates \cite{angelopoulos2023} is now the standard
reference, and it states plainly that the basic guarantee is marginal and that conditional
coverage requires extra machinery. Two score functions matter for classification. The
least-ambiguous-set classifier (LAC) of Sadinle et al.\ \cite{sadinle2019} uses $1-\hat p(y\mid x)$
and yields the smallest sets at a given coverage; adaptive prediction sets \cite{romano2020} trade
size for better conditional behavior. Class-conditional calibration, the Mondrian taxonomy of Vovk
\cite{vovk2013}, splits the calibration set by label and computes a separate quantile for each,
which gives coverage per class at the price of larger sets where a class is hard. Our work does not
invent any of this. We use LAC as the primary score because it is the common default and because
its small sets make the marginal failure starkest; in this revision we add a full comparison
against marginal APS \cite{romano2020} as a second baseline, and we use Mondrian calibration as the
remedy.

\textbf{Conformal prediction in cheminformatics.} Conformal methods entered drug discovery largely
through the work of Norinder, Carlsson, and colleagues, who reframed the classical
applicability-domain question as a coverage question \cite{norinder2014}, and through reviews
aimed at the discovery pipeline \cite{eklund2015}. That line of work has mostly emphasized what
conformal prediction provides---valid sets, an interpretable confidence---rather than where the
default formulation breaks. The imbalance problem is acknowledged in passing in the general
literature but, to our reading, has not been measured systematically on standard molecular
benchmarks, across architectures, or connected to decision cost. That is the space this paper
occupies.

\textbf{Calibration is a separate axis.} It is tempting to assume that a miscalibrated model is the
cause and a calibration fix is the cure. Guo et al.\ \cite{guo2017} showed that modern networks are
often overconfident, and post-hoc scaling helps. We find the conformal failure is only loosely
related to ordinary calibration error: our base models are at worst moderately miscalibrated
(expected calibration error around 0.04--0.10), yet marginal conformal coverage still collapses on
the minority. In this revision we go further and show that even a direct intervention on the
classifier's training objective---class-weighted loss, the standard first response to
imbalance---does not fix the conformal failure either. Calibration, class weighting, and
conditional coverage are three different properties, and fixing the first two does not fix the
third.

\textbf{On the LIT-PCBA benchmark.} A concurrent, independent audit \cite{huang2025} documents
substantial data-leakage and redundancy issues within the curated LIT-PCBA benchmark
\cite{tran2020} itself---duplicate compounds shared across nominally disjoint training and
validation splits, and low-diversity query ligand sets. This is a caution about LIT-PCBA
specifically as a ranking/retrieval benchmark; it does not bear directly on our use here, since we
draw on the raw PubChem primary-screen assay underlying one LIT-PCBA target (AID 651631, TP53) as
an independent classification dataset rather than using LIT-PCBA's curated splits, and we run our
own independent duplicate/label-conflict audit on every dataset we use (Section~3.1.1).

\textbf{Molecular representations.} We deliberately span three representation families. Circular
(ECFP-style) fingerprints \cite{rogers2010} feed a random forest \cite{breiman2001}. A graph
convolutional network \cite{kipf2017} built with PyTorch Geometric \cite{fey2019} learns from the
atom--bond graph. ChemBERTa, a RoBERTa-style transformer pretrained on millions of SMILES strings
\cite{chithrananda2020}, supplies sequence embeddings. These have essentially nothing in common at
the level of how a molecule is encoded, which is what makes their shared failure informative.

\section{Methods}

\subsection{Datasets}

We use five binary classification tasks (Table~\ref{tab:datasets}), chosen to span a wide range of
class imbalance and, in this revision, both ADMET/toxicity endpoints and a genuine bioactivity
high-throughput screen. Blood--brain-barrier penetration (BBBP) \cite{martins2012} labels whether a
compound crosses the barrier; the minority class is the non-penetrant compounds. BACE
\cite{subramanian2016} labels $\beta$-secretase 1 inhibition and is close to balanced. From Tox21
\cite{huang2016} we take the SR-ARE stress-response endpoint. ClinTox provides the CT-TOX endpoint
(failure of clinical trials for toxicity); it is the most imbalanced of the four MoleculeNet
\cite{wu2018} tasks we use and serves as an external check relative to BACE, BBBP, and Tox21.
\textbf{New in this revision:} we add PubChem assay AID~651631, the primary-screen activity outcome
for a p53-pathway high-throughput screen, and the same assay underlying the TP53 target in the
LIT-PCBA benchmark \cite{tran2020}. We use PubChem's own primary-screen Active/Inactive outcome
directly, \emph{not} LIT-PCBA's AVE-debiased, dose-response-confirmed active set; we flag this
distinction explicitly because the two differ substantially in size (602 primary-hit actives here
vs.\ 79 confirmed actives in the curated LIT-PCBA release) and are not interchangeable. We use this
dataset as a genuine, real-world bioactivity-screening case study, independent of any claim about
LIT-PCBA's curation. SMILES strings that RDKit \cite{rdkit} could not parse were dropped; the
counts in Table~\ref{tab:datasets} are post-filtering and post-deduplication (Section~3.1.1), and
are identical across all experiments on a given dataset.

\begin{table}[h]
\centering
\caption{Datasets after parsing, cleaning, and deduplication. ``Minority rate'' is the fraction of
compounds in the smaller class.}
\label{tab:datasets}
\begin{tabular}{lrrrl}
\toprule
Dataset & Compounds & Class 0 & Class 1 & Minority rate \\
\midrule
BACE & 1513 & 822 & 691 & 0.457 \\
BBBP & 1965 & 465 & 1500 & 0.237 \\
Tox21 SR-ARE & 5825 & 4883 & 942 & 0.162 \\
ClinTox CT-TOX & 1440 & 1347 & 93 & 0.065 \\
TP53 (PubChem AID 651631) & 6080 & 5613 & 467 & 0.077 \\
\bottomrule
\end{tabular}
\end{table}

Most experiments use random train/calibration/test splits, which preserve the exchangeability
conformal prediction assumes. For the robustness check in Section~4.8 we instead use scaffold
splits: molecules are grouped by Bemis--Murcko scaffold \cite{bemis1996} and whole groups are
assigned to a single partition, so that test compounds carry skeletons unseen in training and
calibration. Scaffold splits break exchangeability by design and so test behavior under realistic
distribution shift.

\subsubsection{Duplicate and label-conflict audit (new)}

Before splitting, we canonicalized all SMILES with RDKit and checked for duplicate compounds within
each dataset. BACE and Tox21 SR-ARE contained no duplicated canonical SMILES. BBBP contained 64
duplicated canonical SMILES (3.1\% of the pre-dedup dataset), of which 10 groups carried
conflicting labels across duplicate entries. ClinTox contained 19 duplicated canonical SMILES
groups (1.3\% of the pre-dedup dataset), \emph{all} 19 of which were label-conflicting---together
touching 19 of the original 112 minority-class (toxic) compounds, or 17\% of the minority class
underlying the dataset's headline coverage result. We removed all label-conflicting duplicate
groups and collapsed label-consistent duplicates to a single row before any train/calibration/test
split. All results reported in this paper, including Table~\ref{tab:datasets}, are
post-deduplication. The core finding is unchanged by this cleaning step (Table~\ref{tab:main}).

\subsection{Models and representations}

For the main conformal experiments and all cross-dataset work we use a random forest
\cite{breiman2001} (300 trees on BACE, BBBP, ClinTox, and TP53; 150 trees on Tox21 SR-ARE to fit
compute budget, directionally stable across the tree counts tested) on 2048-bit Morgan fingerprints
of radius 2, computed with RDKit's MorganGenerator. We selected the random forest after comparing
three standard classifiers on BBBP under stratified five-fold cross-validation (Section~4.1); it
gave the best AUC and competitive balanced accuracy, and as a bagged ensemble it produces
reasonable probability estimates out of the box.

For the architecture study we add two more representations. The graph network is a two-layer graph
convolutional model \cite{kipf2017} (GCNConv 32$\to$64$\to$64, mean pooling, a linear read-out to
two classes) implemented in PyTorch Geometric \cite{fey2019}, trained for 60 epochs with Adam at
learning rate $10^{-3}$. Atom features are a 32-dimensional one-hot encoding of element, degree,
hydrogen count, hybridization, aromaticity, and formal charge. The transformer is ChemBERTa
(DeepChem/ChemBERTa-77M-MTR) \cite{chithrananda2020}, used as a frozen feature extractor: we
mean-pool its final hidden states over the SMILES tokens to a 384-dimensional embedding and fit a
logistic-regression head. Freezing the transformer isolates the representation from any tuning
choices on our side.

\subsection{Conformal prediction}

We use split (inductive) conformal prediction. Each dataset is partitioned into a proper training
set (50\%), a calibration set (25\%), and a test set (25\%). The base model is fit on the training
split only. For a model output $\hat p(y\mid x)$ we score each calibration point with the LAC
nonconformity score \cite{sadinle2019}
\begin{equation}
s_i = 1 - \hat p(y_i \mid x_i). \label{eq:1}
\end{equation}
For target coverage $1-\alpha$ with $\alpha=0.10$, the marginal threshold is the conformal quantile
\begin{equation}
\hat q = \text{the } \left\lceil (n+1)(1-\alpha) \right\rceil / n \text{ empirical quantile of } \{s_i\}, \label{eq:2}
\end{equation}
where $n$ is the calibration-set size, and the prediction set for a test point is
\begin{equation}
C(x) = \{ y : \hat p(y\mid x) \ge 1-\hat q \}. \label{eq:3}
\end{equation}
The finite-sample correction in Eq.~\ref{eq:2} is what gives the marginal guarantee
$\Pr(y\in C(x)) \ge 1-\alpha$.

Class-conditional (Mondrian) calibration changes one thing. Instead of a single threshold it
computes a separate quantile $\hat q_k$ from the calibration points of each class $k$:
\begin{equation}
\hat q_k = \text{the conformal quantile of } \{s_i : y_i = k\}, \quad
C(x) = \{ y : \hat p(y\mid x) \ge 1-\hat q_y \}. \label{eq:4}
\end{equation}
This enforces $\Pr(y\in C(x)\mid y=k) \ge 1-\alpha$ for every class, which is precisely the
per-class guarantee the marginal version lacks.

\textbf{New in this revision.} We additionally compute marginal calibration under the adaptive
prediction sets (APS) score of Romano, Sesia, and Cand\`es \cite{romano2020}, which sorts class
probabilities in decreasing order and includes labels by cumulative probability mass rather than a
single LAC threshold, to test whether a more modern conformal score avoids the marginal failure
without resorting to class-conditional calibration.

We report coverage overall and within each class, and mean prediction-set size. Conformal
experiments are repeated over 5--10 random train/calibration/test splits per dataset (fewer seeds
on the largest dataset, Tox21 SR-ARE, to fit compute budget); we report the mean and, where
relevant, the standard deviation across seeds. All reported coverage and set-size numbers are
computed on the held-out 25\% test split, never the calibration split.

\subsection{Class-weighted ablation (new)}

To test whether the failure is a symptom of an unweighted classifier rather than of marginal
calibration itself, we repeat the full marginal/Mondrian LAC pipeline with
\texttt{class\_weight=`balanced'} in the random forest on BBBP, Tox21 SR-ARE, and ClinTox, and
compare marginal minority coverage under weighted vs.\ unweighted training.

\subsection{Selective prediction}

A conformal classifier can also be read as an accept/abstain rule. A singleton prediction set is a
confident, actionable call; a two-element set $\{0,1\}$ is an abstention---the model declines and,
in a screening setting, the compound goes to the lab. We report the acceptance rate (fraction of
singletons) and the accuracy on accepted compounds, overall and for the minority class, under both
calibration schemes.

\subsection{Decision-cost model (revised for direction)}

To turn coverage into something a project manager would recognize, we score minority-class
decisions with a simple utility
\begin{equation}
U = B_{tp} N_{tp} - C_{fp} N_{fp} - C_{fn} N_{fn} - C_{lab} N_{sent}, \label{eq:5}
\end{equation}
where $N_{tp}, N_{fp}, N_{fn}$ count correct and incorrect confident calls, $N_{sent}$ counts
abstentions routed to the lab, and $B, C$ are the corresponding benefit and costs, computed over
1000 scaled minority-class decisions. \textbf{Revision:} the original formulation implicitly
treated class~1 as the class to avoid on every dataset. This is correct for Tox21 SR-ARE and
ClinTox (class~1 = toxic, the outcome a screen should catch) but backwards for BACE, where class~1
is the active $\beta$-secretase inhibitor---the actual goal of the screen---so a missed active
(false negative on class~1) is the costly error, not a false positive. We now score utility against
whichever class is both the numerical minority \emph{and} the class a screening program cannot
afford to miss, with an asymmetric cost structure ($C_{fn}=15$ vs.\ $C_{fp}=5$) for BACE and the
toxicity datasets, reflecting that missing the class of interest is more expensive than one extra
lab assay. For BBBP, whose screening goal (finding CNS-penetrant vs.\ peripherally restricted
compounds) is program-dependent, we retain the symmetric cost structure of the original
formulation ($C_{fp}=C_{fn}=5$).

\begin{table}[h]
\centering
\caption{Base classifiers on BBBP under stratified five-fold cross-validation (mean over folds,
deduplicated data). ROC-AUC reported with its standard deviation.}
\label{tab:table2}
\begin{tabular}{lrrr}
\toprule
Classifier & ROC-AUC & Balanced acc. & F1 \\
\midrule
Random forest & 0.924 $\pm$ 0.012 & 0.794 & 0.931 \\
Hist.\ gradient boost & 0.917 $\pm$ 0.007 & 0.813 & 0.926 \\
Logistic regression & 0.913 $\pm$ 0.010 & 0.812 & 0.925 \\
\bottomrule
\end{tabular}
\end{table}

\subsection{Scaffold blind-spot analysis}

To find where the marginal model is confidently wrong, we reduce each compound to its
Bemis--Murcko scaffold \cite{bemis1996}---its ring and linker skeleton, side chains
removed---using RDKit. On the BBBP test split we flag every minority compound whose marginal
prediction set is a confident singleton that excludes the true (minority) label. We then count
these confident errors per scaffold and rank the scaffolds with at least two minority test
compounds.

\subsection{Statistical testing}

For each architecture we compare minority coverage under Mondrian versus marginal calibration with
a one-sided Wilcoxon signed-rank test over the ten paired seeds. We also report expected
calibration error (ECE) \cite{guo2017} and Brier scores for the base models.

\section{Results}

\subsection{Base classifier selection and calibration}

Table~\ref{tab:table2} compares three classifiers on BBBP. The random forest leads on AUC (0.924)
with a histogram gradient-boosting model close behind and slightly ahead on balanced accuracy. We
carried the random forest forward as the base model for all conformal experiments.

The base models are not badly calibrated. Out-of-fold expected calibration error and Brier scores
(Figure~\ref{fig:reliability}) show only mild deviations from the diagonal across all three
classifiers. We make a point of this because it rules out the obvious explanation. The conformal
failure below is not a symptom of a wildly overconfident model. It survives even when ordinary
calibration is fine.

\begin{figure}[h]
\centering
\includegraphics[width=0.95\textwidth]{../figures/figure1_reliability.png}
\caption{Reliability diagrams for three base classifiers on BBBP (out-of-fold, deduplicated data).
Deviation from the diagonal is mild. Ordinary calibration being acceptable is what makes the
conformal coverage failure in Figure~\ref{fig:bbbpcoverage} surprising.}
\label{fig:reliability}
\end{figure}

\subsection{Marginal coverage collapses on the minority class}

At a 90\% target, marginal conformal prediction on BBBP behaves exactly as advertised in aggregate:
overall coverage is on target. Split that number by class and it falls apart. Majority-class
coverage is inflated well above target; minority-class coverage is only 59.2\% (Table~\ref{tab:main}).
The method buys its global guarantee by over-covering the easy class and under-covering the hard
one (Figure~\ref{fig:bbbpcoverage}). A practitioner watching overall coverage would see a number
that looks like the guarantee working. The compounds they care about are covered barely more than
half the time.

\begin{figure}[h]
\centering
\includegraphics[width=0.7\textwidth]{../figures/figure2_bbbp_coverage.png}
\caption{Per-class coverage on BBBP at a 90\% target (random forest, deduplicated data, means over
seeds). Marginal calibration over-covers the majority and leaves the minority far below target.
Mondrian calibration brings both classes to target.}
\label{fig:bbbpcoverage}
\end{figure}

Switching to Mondrian calibration repairs this. Minority coverage rises to target on every dataset
(Table~\ref{tab:main}). The cost is set size: mean prediction-set size grows under Mondrian
calibration on every dataset where the marginal gap was severe. In other words the marginal method
was producing tight sets partly by omitting the true minority label; the conditional method pays
for honest coverage with slightly larger sets on the hard class, which is the correct trade.

\begin{table}[h]
\centering
\caption{Core results across all five datasets (random forest; means over 5--8 seeds).}
\label{tab:main}
\begin{tabular}{lrrrrrl}
\toprule
Dataset & MCC & AUPRC & AUROC & LAC marg.\ min.\ cov.\ & Mondrian min.\ cov.\ & APS marg.\ min.\ cov.\ (set size) \\
\midrule
BACE & 0.62 & 0.85 & 0.89 & 92.3\% & 92.2\% & 99.9\% (1.99) \\
BBBP & 0.64 & 0.96 & 0.90 & 59.2\% & 88.8\% & 100.0\% (2.00) \\
Tox21 SR-ARE & 0.29 & 0.48 & 0.79 & 40.5\% & 92.3\% & 99.9\% (2.00) \\
ClinTox & 0.19 & 0.26 & 0.75 & 2.5\% & 98.3\% & 100.0\% (2.00) \\
TP53 (PubChem) & 0.27 & 0.41 & 0.82 & 2.1\% & 89.1\% & 100.0\% (2.00) \\
\bottomrule
\end{tabular}
\end{table}

\subsection{Severity scales with imbalance}

The same experiment across datasets lines up with their imbalance (Table~\ref{tab:main},
Figure~\ref{fig:crossdataset}). On near-balanced BACE (minority rate 0.457) marginal calibration is
already close to target, and Mondrian leaves it essentially unchanged. As imbalance deepens through
BBBP, Tox21 SR-ARE, ClinTox, and the new TP53 PubChem dataset, marginal minority coverage collapses
further and further while Mondrian calibration holds every dataset near its 90\% target. The
ordering is monotone: the rarer the class, the worse marginal calibration treats it, and BACE shows
the effect switches off when imbalance does.

\begin{figure}[h]
\centering
\includegraphics[width=0.85\textwidth]{../figures/figure3_cross_dataset.png}
\caption{Marginal minority-class coverage against a 90\% target (dashed) across all five datasets,
ordered by imbalance. The shortfall deepens from BACE through BBBP, Tox21, ClinTox, to the new
TP53 PubChem bioactivity dataset; Mondrian calibration returns every dataset to target.}
\label{fig:crossdataset}
\end{figure}

To name the failure with a single number, we use the minority coverage gap,
$\mathrm{MCG} = (1-\alpha) - \mathrm{cov}_{\min}$: how far minority coverage falls below the
promised target. A companion, the coverage imbalance index
$\mathrm{CII} = |\mathrm{cov}_{\mathrm{maj}} - \mathrm{cov}_{\min}|$, captures how unevenly
coverage is split between classes. These are deliberately simple; the value of a diagnostic is that
it is easy to compute and report, not that it is mathematically deep. Under marginal calibration
MCG is largest on the most imbalanced datasets (ClinTox and TP53) and smallest on BACE; under
Mondrian it collapses to within a few points of zero on all five.

We also tested ClinTox and, new in this revision, the PubChem TP53 dataset as external checks; they
appear nowhere else in our experiments except as standalone case studies. Both show near-total
collapse under marginal calibration and near-total recovery under Mondrian calibration, consistent
with the pattern across all five datasets. Section~5 shows why this ordering is not a coincidence.

\subsection{The failure is architecture-agnostic}

\begin{table}[h]
\centering
\caption{Minority-class coverage on BBBP at a 90\% target across three model architectures (ten
seeds each). The failure and its repair appear in all three.}
\label{tab:table4}
\begin{tabular}{lllll}
\toprule
Architecture & Representation & Marginal cov.\ & Mondrian cov.\ & Wilcoxon $p$ \\
\midrule
Random forest & ECFP fingerprints & 64.8\% & 90.3\% & $9.8\times 10^{-4}$ \\
Graph network & atom--bond graph & 63.3\% & 90.9\% & $9.8\times 10^{-4}$ \\
ChemBERTa & SMILES (frozen) & 72.4\% & 90.1\% & $9.8\times 10^{-4}$ \\
\bottomrule
\end{tabular}
\end{table}

If the collapse were specific to fingerprints and random forests, it would be a quirk worth a
footnote. It is not. We repeated the BBBP experiment with two more representations that share
nothing with the first: a graph network reading the molecular graph, and a frozen chemical language
model reading the SMILES string. Both collapse the same way (Table~\ref{tab:table4}). The Wilcoxon
test is at its floor for ten paired seeds in all three cases, with every seed pointing the same
direction.

One detail in Table~\ref{tab:table4} is worth dwelling on. The transformer's marginal coverage is
noticeably higher than the other two. ChemBERTa was pretrained on millions of molecules, so it
arrives with better-calibrated minority predictions, and the marginal threshold therefore mistreats
the minority a little less. But ``a little less'' is still a substantial miss. The vulnerability
shrinks with a stronger pretrained model; it does not disappear. We read this as evidence that the
collapse is driven by the marginal calibration step interacting with baseline minority calibration,
not by anything architectural---a better representation raises the floor without removing the
problem.

\subsection{Aggregate metrics hide the failure}

The selective-prediction view on BBBP shows why a normal evaluation would miss all of this. Under
marginal calibration the model accepts (returns a confident singleton for) the large majority of
compounds at high accuracy on those accepted---an excellent-looking summary. Restrict to the
minority class and the same model's selective accuracy on accepted minority compounds drops
sharply. It is confidently wrong on minority compounds close to half the time, while the headline
accuracy on accepted compounds looks fine. Mondrian calibration changes the behavior in the right
direction: larger sets push more genuinely uncertain minority cases to abstention, and minority
selective accuracy improves substantially. The overall accept/accuracy figures barely move between
the two schemes. That is the trap. Everything you would normally look at says the marginal model is
fine.

\subsection{Does rebalancing the classifier fix it? No (new)}

A natural response to class imbalance is to reweight the classifier's training loss. We repeat the
BBBP, Tox21 SR-ARE, and ClinTox experiments with \texttt{class\_weight=`balanced'} and find marginal
minority coverage essentially unchanged: BBBP 59.2\%$\to$60.7\%, Tox21 SR-ARE 40.5\%$\to$34.7\% (no
improvement), ClinTox 2.5\%$\to$0.0\% (no improvement). Class weighting shifts the classifier's
decision boundary but does not change the shape of the pooled nonconformity-score distribution
enough to move the single marginal quantile off the majority mass (Figure~\ref{fig:nonconformity}).
This is additional evidence for the mechanism in Section~5: the failure lives in the calibration
step, not the base classifier, and no amount of reweighting the classifier substitutes for
calibrating per class.

\subsection{Adaptive prediction sets reach coverage by becoming uninformative (new)}

Both reviewers of the previous submission of this manuscript asked whether more recent conformal
methods avoid this failure without Mondrian calibration. Table~\ref{tab:main} shows marginal APS
reaching near-total minority coverage (99.9--100\%) on every dataset---but mean prediction-set size
under APS is $\approx 2.0$ everywhere, meaning APS achieves its guarantee by returning the trivial
$\{0,1\}$ set for nearly every minority compound. APS therefore fails differently than LAC, not
less: LAC is confidently and often wrong on the minority class, while marginal APS is uninformatively
``correct''---it essentially never commits to a minority-class call, which is operationally
equivalent to a policy of always abstaining and is no more useful for screening decisions than that
policy would be. Class-conditional calibration is the only one of the three schemes tested that
returns both correct minority coverage and prediction sets small enough to act on.

\subsection{The cost of false confidence, direction-corrected}

Feeding the measured minority numbers into the utility of Eq.~\ref{eq:5} makes the stakes concrete.
Table~\ref{tab:utility} reports utility under the corrected, direction-aware cost structure
(Section~3.5). The qualitative conclusion strengthens relative to the original submission: under
asymmetric, direction-corrected costs, Mondrian calibration's advantage is larger everywhere it
matters (Tox21, ClinTox, BBBP), and the two schemes are statistically indistinguishable only on
BACE---the one dataset that is already near-balanced and where the conservation identity predicts
the marginal failure should be negligible. This is the identity's ordering prediction showing up a
second time, independently, in the cost analysis.

\begin{table}[h]
\centering
\caption{Utility over 1000 scaled minority-class decisions, direction-corrected costs.}
\label{tab:utility}
\begin{tabular}{lllrr}
\toprule
Dataset & Screening goal & Cost structure & Marginal $U$ & Mondrian $U$ \\
\midrule
BACE & Find active inhibitors & $C_{fn}\gg C_{fp}$ & $-700$ & $-671$ \\
BBBP & Symmetric (program-dependent) & $C_{fn}=C_{fp}$ & $-1{,}569$ & $+86$ \\
Tox21 SR-ARE & Avoid toxic compounds & $C_{fn}\gg C_{fp}$ & $-9{,}484$ & $-1{,}141$ \\
ClinTox & Avoid toxic compounds & $C_{fn}\gg C_{fp}$ & $-11{,}959$ & $-179$ \\
\bottomrule
\end{tabular}
\end{table}

To make ``robust to the exact constants'' more than an assertion, we swept the cost of a confident
wrong call ($C_{fp}=C_{fn}$) from 1 to 50 and the cost of a lab assay ($C_{lab}$) from 0.2 to 2 on
the BBBP minority-decision framing (Figure~\ref{fig:utilityheatmap}). The Mondrian advantage
$\Delta U = U_{mond} - U_{marg}$ is positive everywhere on the grid, and it grows as confident
errors get more expensive---which is the regime any safety-critical screen lives in. The conclusion
does not depend on our particular choice of cost constants.

\begin{figure}[h]
\centering
\includegraphics[width=0.75\textwidth]{../figures/figure4_utility_heatmap.png}
\caption{Sensitivity of the Mondrian utility advantage $\Delta U = U_{mond}-U_{marg}$ to the cost
assumptions, on BBBP minority decisions ($n=1000$). $\Delta U$ is positive across every combination
of error cost and lab-assay cost, so class-conditional calibration is the better decision rule
regardless of where a campaign sets its costs.}
\label{fig:utilityheatmap}
\end{figure}

\subsection{Where the model is blind: a scaffold case study}

The confident minority errors are not spread evenly across chemical space. They pile up on a few
scaffolds (Figure~\ref{fig:scaffolds}). On the deduplicated BBBP test split (seed 0) the worst
offender is the bare benzene scaffold: of 8 minority compounds reducing to a benzene core, the
marginal model is confidently wrong on 5 (62.5\%). Two other generic scaffolds each show total or
near-total confident misclassification among their (small) minority compound counts. The common
thread is telling. These are simple, generic skeletons that occur in \emph{both} classes. A bare
benzene ring carries almost no information about whether a compound crosses the blood--brain
barrier, so the model defaults to the majority call and does it confidently---and it is exactly the
minority compounds hiding inside that generic chemistry that get burned. Complex, distinctive
scaffolds with several minority test compounds show markedly lower confident-error rates. The blind
spot is not exotic chemistry the model never saw. It is the most ordinary chemistry, where the
minority signal is drowned out. That is a sharper and more uncomfortable statement than ``the model
fails on weird molecules,'' and it is precisely the regime a marginal guarantee papers over.

\begin{figure}[h]
\centering
\includegraphics[width=0.9\textwidth]{../figures/figure5_scaffolds.png}
\caption{The Bemis--Murcko scaffolds with the most confident minority-class errors under marginal
calibration on BBBP (deduplicated data, seed 0). The worst blind spots are generic cores that
appear in both classes and carry little class signal.}
\label{fig:scaffolds}
\end{figure}

\textbf{One sentence connecting this to Mondrian calibration, added per reviewer request:} under
Mondrian calibration, the minority-class threshold widens enough that many of these previously
confidently-wrong compounds on generic scaffolds move from a confident-but-wrong singleton into the
two-label abstention set, which is the mechanism by which Mondrian calibration converts confident
errors into flagged uncertainty rather than eliminating the underlying difficulty of the chemistry.

\subsection{Robustness to realistic scaffold splits}

Random train/test splits let near-identical molecules land on both sides, which flatters any model
and, in principle, could flatter the conformal calibration too. The honest test is a scaffold
split, where whole Bemis--Murcko scaffolds are assigned to a single partition so that test compounds
carry skeletons absent from training and calibration. This deliberately violates the exchangeability
that conformal coverage assumes, so it probes robustness under distribution shift rather than a
clean guarantee. We repeated the BBBP experiment with ten random scaffold splits. The collapse does
not go away; it deepens. Marginal minority coverage falls further than under random splits, and
Mondrian recovers it, over-covering slightly rather than landing exactly on target, which is the
expected and safe direction when exchangeability is broken.

\textbf{Rationale, added per reviewer request:} scaffold splitting deliberately breaks the
exchangeability conformal prediction requires, so coverage guarantees under scaffold splits are not
formally exact in the way random-split guarantees are. We include this experiment not as a second
valid coverage claim but as a stress test: real deployment means scoring molecules with scaffolds
the model has never calibrated on, and if the marginal-vs-Mondrian gap only appeared under the
idealized random-split assumption, that would substantially limit the practical relevance of our
findings. The gap does not shrink under scaffold splits---if anything it deepens
slightly---which is the relevant robustness check for a practitioner deciding whether to trust
conformal prediction on genuinely novel chemistry, even though neither scheme's numeric guarantee
is formally exact in this regime.

\section{Mechanism: why marginal calibration under-covers the minority}

The experiments establish that the failure is real, general, and costly. They do not yet say why.
Here we give a short argument that explains the effect, predicts its size from quantities known
before any model is trained, and accounts for the ordering across datasets.

\textbf{Conservation of coverage.} Split conformal prediction guarantees only the \emph{overall}
coverage, $\Pr(y\in C(x)) \ge 1-\alpha$. Overall coverage is the class-weighted average of the
per-class coverages, with weights equal to the class prevalences $\pi_{\mathrm{maj}}, \pi_{\min}$:
\begin{equation}
1-\alpha = \pi_{\mathrm{maj}}\,\mathrm{cov}_{\mathrm{maj}} + \pi_{\min}\,\mathrm{cov}_{\min}. \label{eq:6}
\end{equation}

\begin{figure}[h]
\centering
\includegraphics[width=0.75\textwidth]{../figures/figure6_nonconformity.png}
\caption{Calibration nonconformity scores on BBBP, split by class (deduplicated data, seed 0). The
majority (penetrant) scores concentrate near zero; the minority (non-penetrant) scores spread to
higher values. The single marginal threshold is pulled toward the majority mass and slices through
the minority distribution, excluding much of its upper tail---the source of the under-coverage.
Mondrian gives each class its own threshold; the minority threshold sits far enough right to cover
its tail.}
\label{fig:nonconformity}
\end{figure}

Solving for the minority coverage and rearranging gives the relation that drives the whole
phenomenon:
\begin{equation}
(1-\alpha) - \mathrm{cov}_{\min} = \frac{\pi_{\mathrm{maj}}}{\pi_{\min}}\left(\mathrm{cov}_{\mathrm{maj}} - (1-\alpha)\right). \label{eq:7}
\end{equation}
\begin{center}
\emph{minority shortfall} \hspace{2.5cm} \emph{majority surplus}
\end{center}

Coverage is conserved: whatever the majority is given above target, the minority must lose below
it. And the loss is not one-for-one---it is amplified by the imbalance ratio
$\pi_{\mathrm{maj}}/\pi_{\min}$. At three-to-one imbalance, each point of majority over-coverage
costs three points of minority coverage. The mechanism is visible directly in the calibration
scores (Figure~\ref{fig:nonconformity}): because the model classifies the abundant majority easily,
its nonconformity scores pile up near zero, so the pooled $(1-\alpha)$ quantile---the single
marginal threshold---is dragged down to a value that comfortably covers the majority while failing
to reach the minority's much heavier upper tail. The majority's surplus is exactly the
over-coverage that this low threshold produces.

\textbf{The identity predicts the failure.} Equation~\ref{eq:7} is not a loose analogy; it holds
numerically to within a few points across every dataset we study, and the amplification factor
$\pi_{\mathrm{maj}}/\pi_{\min}$ orders the severity of the collapse across datasets exactly as
measured. Nothing about this is specific to chemistry: the only ingredients are class imbalance and
a model that handles the common class better than the rare one, which is the normal situation in
screening.

\subsection{A validated pre-flight coverage-gap diagnostic (new)}

Equation~\ref{eq:7} \emph{explains} the measured minority coverage gap after marginal calibration
has already been run and $\mathrm{cov}_{\mathrm{maj}}$ measured. It is diagnostic, not predictive
in the useful sense---you need the failure to already exist, and to have already run the full
calibration pipeline, to explain it. We extend it into a pre-flight predictor: given only (a) class
prevalence, known before any modeling, and (b) one cheap marginal-calibration pass to obtain
$\mathrm{cov}_{\mathrm{maj}}$, a practitioner can predict $\mathrm{cov}_{\min}$ before deciding
whether class-conditional calibration is worth implementing. Because coverage is bounded in
$[0,1]$ while the identity is linear and unbounded, we clip the prediction to $[0,1]$---and note
explicitly that the unclipped identity can produce nonsensical negative predictions at high
imbalance (e.g.\ $-40.6\%$ on ClinTox, amplification factor $16.1\times$ in that particular test
split), which is itself informative: it flags that imbalance has passed the point where marginal
calibration can be trusted at all.

\begin{table}[h]
\centering
\caption{Pre-flight diagnostic validation across five datasets spanning 1.2$\times$--16.1$\times$
imbalance.}
\label{tab:diagnostic}
\begin{tabular}{lrrrr}
\toprule
Dataset & Amplification factor & Predicted min.\ cov.\ (clipped) & Actual min.\ cov.\ & Abs.\ error \\
\midrule
BACE & 1.22$\times$ & 88.9\% & 92.3\% & 3.4 pts \\
BBBP & 3.21$\times$ & 62.3\% & 59.2\% & 3.1 pts \\
Tox21 SR-ARE & 5.10$\times$ & 39.8\% & 40.5\% & 0.7 pts \\
ClinTox & 16.14$\times$ & 0.0\% & 2.5\% & 2.5 pts \\
TP53 (PubChem) & 12.01$\times$ & 0.0\% & 2.1\% & 2.1 pts \\
\bottomrule
\end{tabular}
\end{table}

Validated against all five datasets (Table~\ref{tab:diagnostic}), imbalance ratios
1.2$\times$--16.1$\times$, the clipped diagnostic predicts minority coverage to within a mean
absolute error of 2.4 percentage points---including correctly flagging near-total collapse at the
two most extreme imbalance ratios once the identity is bounded. We propose this as a one-line
pre-registration-style check practitioners can run before committing to a full calibration study,
and release it as a single function in the accompanying repository.

\textbf{Limitation, stated explicitly.} This diagnostic is a bounded, first-order approximation,
not an exact law. It requires one calibration pass, not zero, degrades in interpretability above
roughly $10\times$ imbalance where the unclipped identity leaves the $[0,1]$ range entirely, and
has been validated on five datasets---a wider sweep across endpoints and imbalance ratios (noted in
Limitations) would be needed to fit it as a general predictive law rather than a validated trend.

\textbf{Why Mondrian works, and why temperature scaling would not.} Mondrian calibration breaks
Eq.~\ref{eq:6} by refusing to pool. With a separate quantile per class, each class gets its own
$(1-\alpha)$ threshold and the amplification channel is closed---there is no shared budget to
redistribute. This also clarifies why ordinary recalibration is the wrong tool. Rescaling
probabilities (for instance, temperature scaling) shifts every score but cannot change the fact
that a single pooled quantile is dominated by the majority; it adjusts confidence, not the
conditional coverage that Eq.~\ref{eq:7} governs. The fix has to act on the calibration partition,
not on the probabilities. Similarly---and this is new in this revision---reweighting the
classifier's training objective does not act on the calibration partition either, which is why the
class-weighted ablation in Section~4.6 shows no improvement.

\textbf{An illusion of safety.} Put the pieces together and a pattern emerges that is worse than any
single number. A team adopts conformal prediction precisely to stop trusting a model blindly. They
pick $\alpha=0.10$, confirm overall coverage near 90\%, see high accuracy on confident calls, and
conclude the system is safe to act on. Every one of those checks passes while the model covers the
minority class barely more than half the time on some datasets and almost never on others, abstains
on the wrong compounds, and is most confidently wrong on the most ordinary molecules in the library.
The marginal guarantee did not lie---it delivered exactly the average coverage it
promised---but the way it is naturally read, as a per-case assurance, is an illusion of safety on
imbalanced data.

\textbf{What is and is not new here.} The mathematics of conformal prediction is not ours and we
want to be exact about that. The gap between marginal and conditional coverage is foundational
conformal theory \cite{vovk2013,angelopoulos2023}, and Mondrian calibration is the standard answer
to it. Our contribution is empirical, explanatory, and---new in this revision---predictive: that
this known gap is large, common, and invisible under the imbalance that defines virtual screening;
that it holds across three unrelated architectures and a genuine bioactivity screen, with severity
set by baseline minority calibration rather than architecture; that a one-line conservation
identity predicts its size and ordering from class prevalences alone, and can be turned into a
bounded pre-flight check validated across five datasets and a thirteen-fold imbalance range; that
ordinary metrics conceal it; that ordinary class rebalancing does not fix it; that a more modern
conformal score (APS) fails differently rather than not at all; and that the failure localizes to
interpretable chemical substructures with a measurable, direction-aware, cost-robust penalty. The
conservation argument is elementary, and we make no claim otherwise---its worth is that it turns an
observed failure into a predictable one.

\textbf{Practical recommendation.} The operational takeaway is short. On any imbalanced screening
task, report coverage per class, not just overall; run the pre-flight diagnostic before committing
to a full study; and calibrate per class with Mondrian conformal prediction by default rather than
reaching for class-weighted training or a fancier marginal score, neither of which fixes the
problem. The extra cost of Mondrian calibration is a modest increase in prediction-set size,
concentrated on the hard class, which is the honest place to spend it. Treating larger minority
sets as a defect is backwards; they are the method telling the truth about a class it cannot yet
call confidently.

\section{Discussion and Limitations}

Several limitations remain. We study five datasets and one imbalance axis; a fuller sweep across
endpoints and ratios would turn the amplification relation from a validated trend into a fitted
law. Our scaffold-split test uses random scaffold partitions over ten seeds rather than a single
canonical split, which is why its coverage estimates are noisier; the direction is nonetheless
unambiguous. ClinTox and the new TP53 PubChem dataset are both fairly imbalanced with modest
absolute minority-class counts, so their percentages are somewhat coarse at the level of individual
test splits. We froze ChemBERTa rather than fine-tuning it, so its numbers reflect pretrained
representations, not a tuned model. We used a single primary nonconformity score (LAC), though this
revision adds a full comparison against adaptive prediction sets (APS) as a second baseline; a
head-to-head of LAC against APS under \emph{class-conditional} (rather than only marginal)
calibration remains a natural extension. The direction-aware cost model is illustrative---we show
the sign and robustness of the effect, not real campaign economics. The scaffold blind-spot
analysis is run on one seed for interpretability; the identity of the worst scaffolds may shift
across splits even if the generic-core pattern does not. The pre-flight diagnostic is a bounded
first-order approximation, explicitly flagged in Section~5.1 as degrading in interpretability above
roughly tenfold imbalance. Finally, the TP53 PubChem dataset is a primary-screen outcome, not
LIT-PCBA's curated, dose-response-confirmed active set; we use it as an independent bioactivity
case study, not as a claim of alignment with LIT-PCBA's specific curation protocol, and note the
concurrent audit of LIT-PCBA's own data integrity \cite{huang2025} as context for why we did not
simply adopt that benchmark's splits wholesale.

\textbf{Future work.} The natural extensions are a wider dataset sweep to fit the amplification
relation directly and sharpen the pre-flight diagnostic; a head-to-head of LAC against adaptive
prediction sets under \emph{class-conditional} calibration specifically; group-conditional
calibration along axes beyond the label---chemical series, assay source---where silent
under-coverage could hide for the same reason; and a fuller accounting of when class-weighted
training might interact with calibration in ways our three-dataset ablation did not probe.

\section{Conclusion}

Conformal prediction gives virtual screening something it badly needs: a distribution-free
statement about reliability. Read carelessly on imbalanced data, that statement is also a way to be
confidently wrong about the compounds that matter most. We measured marginal conformal prediction
under-covering the minority class by tens of points on ADMET/toxicity endpoints and by nearly
ninety points on a real bioactivity screen, while every aggregate metric stayed healthy; showed the
failure across fingerprint, graph, and transformer models at $p<0.001$ each; showed it survives
class-weighted training and is not solved by a more modern conformal score; turned the paper's own
explanatory identity into a validated pre-flight check; traced the failure to a handful of generic
scaffolds; and priced it with a direction-aware cost model. The remedy is already on the
shelf---calibrate per class. The point of the paper is to make the case that, for screening,
per-class calibration should be the default rather than a refinement, and to give practitioners a
concrete way to see the failure coming before it costs them.

\section*{Data and Code Availability}
MoleculeNet datasets (BBBP, BACE, Tox21, ClinTox) \cite{wu2018} are public. The TP53 dataset is
PubChem assay AID~651631, publicly available via PubChem. Code to reproduce every experiment,
table, and figure, with fixed random seeds, is available at
\url{https://github.com/Muhammadjon-crypto/Conformal-Minority-Screen}. A Zenodo-archived release
with a permanent DOI will be provided on acceptance.

\section*{Funding}
This research did not receive any specific grant from funding agencies in the public, commercial,
or not-for-profit sectors.

\section*{Author Contributions}
Muhammadjon Tursunbadalov and Mustafojon Tursunbadalov jointly designed the study, implemented the
experiments, analyzed the results, and wrote the manuscript.

\section*{Conflicts of Interest}
The authors declare no competing interests.

\begin{thebibliography}{99}

\bibitem{vovk2005} Vovk, V.; Gammerman, A.; Shafer, G. \emph{Algorithmic Learning in a Random
World}; Springer: New York, 2005.

\bibitem{angelopoulos2023} Angelopoulos, A. N.; Bates, S. A Gentle Introduction to Conformal
Prediction and Distribution-Free Uncertainty Quantification. \emph{Foundations and Trends in
Machine Learning} \textbf{2023}, \emph{16}(4), 494--591.

\bibitem{sadinle2019} Sadinle, M.; Lei, J.; Wasserman, L. Least Ambiguous Set-Valued Classifiers
With Bounded Error Levels. \emph{Journal of the American Statistical Association} \textbf{2019},
\emph{114}(525), 223--234.

\bibitem{romano2020} Romano, Y.; Sesia, M.; Cand\`es, E. Classification with Valid and Adaptive
Coverage. In \emph{Advances in Neural Information Processing Systems (NeurIPS)}, 2020.

\bibitem{vovk2013} Vovk, V. Conditional Validity of Inductive Conformal Predictors. \emph{Machine
Learning} \textbf{2013}, \emph{92}, 349--376.

\bibitem{wu2018} Wu, Z.; Ramsundar, B.; Feinberg, E. N.; Gomes, J.; Geniesse, C.; Pappu, A. S.;
Leswing, K.; Pande, V. MoleculeNet: A Benchmark for Molecular Machine Learning. \emph{Chemical
Science} \textbf{2018}, \emph{9}, 513--530.

\bibitem{rogers2010} Rogers, D.; Hahn, M. Extended-Connectivity Fingerprints. \emph{Journal of
Chemical Information and Modeling} \textbf{2010}, \emph{50}, 742--754.

\bibitem{breiman2001} Breiman, L. Random Forests. \emph{Machine Learning} \textbf{2001}, \emph{45},
5--32.

\bibitem{kipf2017} Kipf, T. N.; Welling, M. Semi-Supervised Classification with Graph Convolutional
Networks. In \emph{Proceedings of the International Conference on Learning Representations
(ICLR)}, 2017.

\bibitem{fey2019} Fey, M.; Lenssen, J. E. Fast Graph Representation Learning with PyTorch
Geometric. \emph{International Conference on Learning Representations (ICLR) Workshop}, 2019.

\bibitem{chithrananda2020} Chithrananda, S.; Grand, G.; Ramsundar, B. ChemBERTa: Large-Scale
Self-Supervised Pretraining for Molecular Property Prediction. \emph{arXiv:2010.09885}, 2020.

\bibitem{guo2017} Guo, C.; Pleiss, G.; Sun, Y.; Weinberger, K. Q. On Calibration of Modern Neural
Networks. In \emph{International Conference on Machine Learning (ICML)}, 2017.

\bibitem{norinder2014} Norinder, U.; Carlsson, L.; Boyer, S.; Eklund, M. Introducing Conformal
Prediction in Predictive Modeling. A Transparent and Flexible Alternative to Applicability Domain
Determination. \emph{Journal of Chemical Information and Modeling} \textbf{2014}, \emph{54},
1596--1603.

\bibitem{eklund2015} Eklund, M.; Norinder, U.; Boyer, S.; Carlsson, L. The Application of Conformal
Prediction to the Drug Discovery Process. \emph{Annals of Mathematics and Artificial Intelligence}
\textbf{2015}, \emph{74}, 117--132.

\bibitem{martins2012} Martins, I. F.; Teixeira, A. L.; Pinheiro, L.; Falcao, A. O. A Bayesian
Approach to in Silico Blood--Brain Barrier Penetration Modeling. \emph{Journal of Chemical
Information and Modeling} \textbf{2012}, \emph{52}, 1686--1697.

\bibitem{subramanian2016} Subramanian, G.; Ramsundar, B.; Pande, V.; Denny, R. A. Computational
Modeling of $\beta$-Secretase 1 (BACE-1) Inhibitors Using Ligand-Based Approaches. \emph{Journal of
Chemical Information and Modeling} \textbf{2016}, \emph{56}, 1936--1949.

\bibitem{huang2016} Huang, R.; Xia, M.; Nguyen, D.-T.; Zhao, T.; Sakamuru, S.; Zhao, J.; Shahane, S.
A.; Rossoshek, A.; Simeonov, A. Tox21 Challenge to Build Predictive Models of Nuclear Receptor and
Stress Response Pathways as Mediated by Exposure to Environmental Chemicals and Drugs.
\emph{Frontiers in Environmental Science} \textbf{2016}, \emph{3}, 85.

\bibitem{rdkit} Landrum, G. RDKit: Open-Source Cheminformatics; \url{https://www.rdkit.org}
(accessed 2026).

\bibitem{bemis1996} Bemis, G. W.; Murcko, M. A. The Properties of Known Drugs. 1. Molecular
Frameworks. \emph{Journal of Medicinal Chemistry} \textbf{1996}, \emph{39}, 2887--2893.

\bibitem{tran2020} Tran-Nguyen, V.-K.; Jacquemard, C.; Rognan, D. LIT-PCBA: An Unbiased Data Set
for Machine Learning and Virtual Screening. \emph{Journal of Chemical Information and Modeling}
\textbf{2020}, \emph{60}(9), 4263--4273.

\bibitem{huang2025} Huang, A.; Knight, I. S.; Naprienko, S. Data Leakage and Redundancy in the
LIT-PCBA Benchmark. \emph{arXiv:2507.21404}, 2025.

\end{thebibliography}

\end{document}
